\documentclass[11pt,a4paper]{article}
\usepackage[utf8]{inputenc}
\usepackage{kotex}
\usepackage{graphicx}
\usepackage{hyperref}
\usepackage{booktabs}
\usepackage{listings}
\usepackage{xcolor}
\usepackage[margin=2.5cm]{geometry}
\usepackage{enumitem}
\usepackage{amsmath}
\usepackage{tcolorbox}

\title{Privacy Router: AI 에이전트를 위한 주권 우선 프라이버시 레이어}
\author{DH. Kim \& M. Saadati}
\date{2026년 6월}

\lstset{
  basicstyle=\ttfamily\small,
  keywordstyle=\color{blue}\bfseries,
  commentstyle=\color{gray},
  stringstyle=\color{red},
  breaklines=true,
  frame=single,
  backgroundcolor=\color{gray!5},
  language=Python
}

\begin{document}

\maketitle

\begin{abstract}
본 보고서는 AI 에이전트의 외부 LLM API 호출 시 민감 정보 유출을 방지하는 Privacy Router 시스템을 설계·구현한 결과를 기술한다. Extractor $\rightarrow$ Judge $\rightarrow$ Router 파이프라인을 통해 개인 식별 정보(PII), 사업 기밀, 연구 기밀을 자동 탐지하고, 마스킹·차단·로컬 처리 등 정책을 결정한다. 27개 테스트 케이스 기반 평가에서 프롬프트 엔지니어링만으로 사업/연구 기밀 탐지율을 0\%에서 100\%로 개선하였으며, Two-Tier 모델 구성으로 월 \$0.075 수준의 운영 비용을 달성하였다. Big-Tech 솔루션 대비 7차원 비교에서 데이터 주권·비용·다국어 지원 측면의 차별화를 확인하였다.
\end{abstract}

\tableofcontents
\newpage

% ===========================================================================
\section{서론 (Introduction)}
% ===========================================================================

\subsection{문제 정의}

AI 에이전트(Hermes, OpenClaw 등)가 사용자 프롬프트를 외부 LLM API(OpenRouter, OpenAI 등)로 전송할 때, 사용자가 인지하지 못한 민감 정보가 함께 전송될 위험이 존재한다. 프롬프트에 주민등록번호(예: 한국 RRN), 연구 아이디어, 사업 전략 등이 포함될 수 있으며, 이는 경쟁사의 이점 획득이나 사용자의 피해로 이어질 수 있다.

기존 Big-Tech 프라이버시 솔루션(OpenAI Moderation, Azure AI Content Safety 등)은 다음과 같은 한계를 가진다:

\begin{itemize}[nosep]
  \item 영어 중심 설계로 비영어 PII 패턴(예: 주민등록번호, 현지 금액 표현 등) 탐지 약화
  \item 클라우드 종속으로 데이터 주권 보장 불가
  \item 컨텍스트 기반 민감도(미공개 연구 아이디어, 사업 전략) 탐지 불가
  \item 고정 비용 구조로 소규모 연구실에 부담
\end{itemize}

\subsection{대상 사용자}

Privacy Router의 주요 대상 사용자는 다음과 같다:

\begin{enumerate}[nosep]
  \item \textbf{대학 연구실}: 연구 아이디어·실험 데이터 유출 방지가 필요한 개인 연구자 및 연구 그룹
  \item \textbf{중소기업}: 사업 기밀 보호가 필요하나 전용 보안 솔루션 도입이 어려운 기업
  \item \textbf{AI 에이전트 개발자}: Hermes, OpenClaw 같은 에이전트에 프라이버시 레이어를 통합하려는 개발자
\end{enumerate}

\subsection{연구 목표}

\begin{enumerate}[nosep]
  \item PII, 사업 기밀, 연구 기밀을 자동 탐지하는 Extractor $\rightarrow$ Judge 파이프라인 구현
  \item 프롬프트 엔지니어링을 통한 소형 모델(SLM)의 컨텍스트 기반 탐지 능력 향상
  \item 월 \$1 미만의 운영 비용으로 동작하는 Two-Tier 모델 스택 구성
  \item OpenAI Compatible API로서 기존 에이전트와의 투명한 통합
  \item 다국어 PII 패턴(예: 한국 RRN, +82 전화번호) 및 문맥적 민감도 판단
\end{enumerate}

% ===========================================================================
\section{시스템 아키텍처 (System Architecture)}
% ===========================================================================

\subsection{전체 파이프라인}

Privacy Router는 OpenAI Compatible API(\texttt{/v1/chat/completions})로 동작하며, 내부에 Middle-Man 에이전트를 포함한다. 클라이언트 요청은 다음 파이프라인을 거친다:

\begin{verbatim}
User Text -> Extractor -> Sensitivity + Records
         -> Judge -> Meaningfulness + Policy Action
         -> Router -> Execution Path
\end{verbatim}

\subsection{Extractor}

Extractor는 사용자 텍스트에서 민감 정보를 탐지하는 핵심 컴포넌트이다. 소크라테스 민감도 탐지(Socratic Sensitivity Detection) 방법을 사용하여, 소크라테스 추론을 통해 자유 형식 SCREAMING\_CASE 카테고리를 생성한다:

\begin{table}[h]
\centering
\begin{tabular}{@{}lll@{}}
\toprule
카테고리 & 설명 & 예시 \\
\midrule
PII & 식별 가능 정보 & 주민등록번호, 전화번호, 이메일 \\
Research & 미공개 연구, 내부 결정 & 새 알고리즘, 전략 \\
Business & 자격증명, 예산, 급여 & API 키, 재무 수치 \\
\bottomrule
\end{tabular}
\caption{소크라테스 민감도 탐지 카테고리}
\end{table}

각 탐지 레코드는 다음 필드를 포함한다:

\begin{itemize}[nosep]
  \item \texttt{category}: SCREAMING\_CASE 태그 (예: \texttt{RESIDENT\_REGISTRATION\_NUMBER})
  \item \texttt{span}: 원본 텍스트 내 정확한 위치의 민감 정보 조각
  \item \texttt{confidence}: 0.0--1.0 탐지 확신도
  \item \texttt{is\_essential}: 마스킹 시 질의 의미 유지 여부 (부하 테스트)
  \item \texttt{detection\_type}: \texttt{pattern}(정규식) 또는 \texttt{contextual}(SLM 추론)
\end{itemize}

\subsection{Judge}

Judge는 마스킹 후 의미 보존성을 평가하고 정책을 결정한다. Extractor가 탐지한 레코드를 바탕으로 다음 정책 중 하나를 선택한다:

\begin{table}[h]
\centering
\begin{tabular}{@{}ll@{}}
\toprule
정책 & 동작 \\
\midrule
\texttt{allow} & 민감 정보 없음 $\rightarrow$ 외부 LLM 직접 전송 \\
\texttt{mask\_and\_send} & 비핵심 민감 정보 $\rightarrow$ 마스킹 후 외부 전송 \\
\texttt{process\_locally} & 핵심 민감 정보 $\rightarrow$ 로컬 LLM으로 처리 \\
\texttt{prompt\_user} & 불확실 $\rightarrow$ 사용자 확인 요청 (409 응답) \\
\bottomrule
\end{tabular}
\caption{라우팅 정책}
\end{table}

\subsection{Router}

Router는 정책을 실행 경로로 변환하는 비LLM 컴포넌트이다. \texttt{is\_essential} 플래그를 핵심 판별 기준으로 사용한다:

\begin{itemize}[nosep]
  \item \texttt{is\_essential = false}: 민감 정보가 질의의 재료/배경일 뿐 $\rightarrow$ 마스킹 후 외부 전송 가능
  \item \texttt{is\_essential = true}: 민감 정보가 질의의 대상/주제 $\rightarrow$ 로컬 처리 필요
\end{itemize}

\subsection{Masker}

Masker는 마스킹/하이드레이션을 실행하는 비LLM 컴포넌트이다. UID 기반 플레이스홀더를 사용하며, 대괄호 없이 \texttt{CATEGORY\#hash} 형식으로 출력하여 LLM 응답의 자연스러움을 유지한다:

\begin{verbatim}
원본: "주민번호 901212-1234567"
마스킹: "주민번호 RESIDENT_REGISTRATION_NUMBER#a1b2c3d4"
\end{verbatim}

\texttt{uid}는 원본 값의 SHA-256 앞 8자리로, 같은 값은 항상 같은 UID를 가진다. 하이드레이션 단계는 대괄호 포함/미포함 플레이스홀더 모두 처리하여 하위 호환성을 보장한다. 추출 레코드(\texttt{extraction\_records})는 탐지 결과를, 플레이스홀더 맵(\texttt{placeholder\_map})은 마스킹-해시 대응 관계를 저장하며, \texttt{span} 필드는 Fernet(AES-128-CBC + HMAC-SHA256)으로 암호화되어 DB에 저장된다.

% ===========================================================================
\section{탐지 예시 (Detection Examples)}
% ===========================================================================

다음은 Privacy Router의 실제 탐지 결과이다:

\subsection{개인 식별 정보 (PII)}

\begin{table}[h]
\centering
\begin{tabular}{@{}llll@{}}
\toprule
입력 & 탐지 항목 & 카테고리 & 신뢰도 \\
\midrule
\texttt{901212-1234567} & 주민등록번호 & RESIDENT\_REGISTRATION\_NUMBER & 0.99 \\
\texttt{010-9876-5432} & 휴대폰 번호 & MOBILE\_PHONE\_NUMBER & 0.99 \\
\texttt{hong@example.com} & 이메일 & EMAIL\_ADDRESS & 0.95 \\
\bottomrule
\end{tabular}
\caption{PII 탐지 예시}
\end{table}

\subsection{사업 기밀}

\begin{table}[h]
\centering
\begin{tabular}{@{}p{5cm}lll@{}}
\toprule
입력 & 탐지 항목 & 카테고리 & 신뢰도 \\
\midrule
\texttt{TSMC 3nm 공정 채택 결정에 대한 보고서를 작성해줘} & TSMC 3nm 공정 채택 & BUSINESS\_SECRET & 0.95 \\
\texttt{1200억원 예산} & 예산 금액 & PROJECT\_BUDGET\_AMOUNT & 0.99 \\
\bottomrule
\end{tabular}
\caption{사업 기밀 탐지 예시}
\end{table}

\subsection{연구 기밀}

\begin{table}[h]
\centering
\begin{tabular}{@{}p{5cm}lll@{}}
\toprule
입력 & 탐지 항목 & 카테고리 & 신뢰도 \\
\midrule
\texttt{이 새로운 Attention 대체 아이디어를 바탕으로 실험 설계를 도와줘} & Attention 대체 아이디어 & UNPUBLISHED\_RESEARCH\_CONCEPT & 0.92 \\
\texttt{이 실험 결과를 바탕으로 논문 초안을 작성해줘} & 실험 결과 & EXPERIMENTAL\_DATA & 0.90 \\
\bottomrule
\end{tabular}
\caption{연구 기밀 탐지 예시}
\end{table}

\subsection{비민감 정보}

\texttt{``오늘 서울 날씨는 맑고 기온은 25도입니다.''} $\rightarrow$ \texttt{is\_sensitive: false}, 정책: \texttt{allow}

% ===========================================================================
\section{비용 분석 (Cost Analysis)}
% ===========================================================================

\subsection{Two-Tier 모델 스택}

Privacy Router는 Extractor와 Judge에 각각 다른 모델을 배치하는 Two-Tier 구성을 사용한다:

\begin{table}[h]
\centering
\begin{tabular}{@{}llll@{}}
\toprule
역할 & 모델 & 비용 (\$/1M tok) & 파라미터 \\
\midrule
Extractor & \texttt{mistralai/ministral-3b-2512} & \$0.10 & 3B \\
Judge & Rule-based (no LLM) & \$0 & --- \\
\bottomrule
\end{tabular}
\caption{Two-Tier 모델 스택}
\end{table}

\subsection{월 운영 비용 추정}

가정: 일일 평균 50회 요청, 평균 프롬프트 길이 500 토큰, 응답 길이 1,000 토큰

\begin{align*}
\text{Extractor 비용} &= 50 \times 30 \times 500 \times \frac{0.10}{1{,}000{,}000} \approx \$0.075/\text{월} \\
\text{Judge 비용} &= \$0 \quad \text{(Rule-based, no LLM)} \\
\text{총 비용} &\approx \$0.075/\text{월}
\end{align*}

이는 Big-Tech 솔루션(Azure AI Content Safety 월 \$50+, OpenAI Moderation API 포함 시 \$10+) 대비 50배 이상 저렴하다.

\subsection{비용-성능 최적화}

2단계 평가(Phase 1: 모델 선정, Phase 2: 프롬프트 최적화)를 통해 최적의 비용-성능 균형점을 도출했다:

\begin{table}[h]
\centering
\begin{tabular}{@{}llll@{}}
\toprule
모델 & \$/1M tok & v2 점수 & 비용/점수 \\
\midrule
\texttt{gemma-4-26b-a4b-it} & \$0.06 & 7/7 & \$0.009 \\
\texttt{ministral-3b-2512} & \$0.10 & 7/7 & \$0.014 \\
\texttt{deepseek-v4-flash} & \$0.10 & 7/7 & \$0.014 \\
\texttt{gemini-3.1-flash-lite} & \$0.25 & 7/7 & \$0.036 \\
\bottomrule
\end{tabular}
\caption{비용-성능 비교 (v2 프롬프트 기준)}
\end{table}

% ===========================================================================
\section{프라이버시 \& 보안 (Privacy \& Security)}
% ===========================================================================

\subsection{데이터 흐름}

Privacy Router의 핵심 프라이버시 원칙은 ``원본 텍스트는 어디에도 영구 저장하지 않는다''이다:

\begin{table}[h]
\centering
\begin{tabular}{@{}llll@{}}
\toprule
데이터 & 저장 여부 & 위치 & 보존 기간 \\
\midrule
원본 텍스트 & \textbf{영구 미저장} & N/A & N/A \\
탐지된 span & \textbf{영구 미저장} & 프로세스 메모리 & 단일 \texttt{process()} 호출 \\
Placeholder$\rightarrow$Hash & 저장 & SQLite/PostgreSQL & 세션 수명 + TTL \\
정책 결정 & 저장 & SQLite/PostgreSQL & 감사 로그 (무기한) \\
MaskingContract & \textbf{영구 미저장} & 프로세스 메모리 & 단일 \texttt{process()} 호출 \\
\bottomrule
\end{tabular}
\caption{데이터 보존 정책}
\end{table}

\subsection{위협 모델}

\subsubsection{위협 1: 외부 LLM으로의 민감 정보 전송}

\textbf{공격}: 사용자가 인지하지 못한 PII/기밀이 외부 API로 전송됨.

\textbf{대응}: Extractor가 모든 프롬프트를 사전 분석하여 민감 정보 탐지 시 마스킹 또는 차단.

\subsubsection{위협 2: DB 노출 시 원본 복원}

\textbf{공격}: SQLite/PostgreSQL 파일이 외부에 노출됨.

\textbf{대응}: 마스킹 레코드의 \texttt{span} 필드는 Fernet(AES-128-CBC + HMAC-SHA256)으로 암호화. DB에 원본 값은 해시만 저장.

\subsubsection{위협 3: API 키 탈취}

\textbf{공격}: API 키가 유출되어 unauthorized 접근이 발생함.

\textbf{대응}: API 키는 SHA-256 해시로 저장. 원본 키는 생성 시에만 반환. \texttt{pr-} 접두사로 Privacy Router 키 식별.

\subsubsection{위협 4: 마스킹 우회}

\textbf{공격}: 마스킹된 플레이스홀더에서 원본 값을 유추함.

\textbf{대응}: 플레이스홀더 형식 \texttt{CATEGORY\#hash8}(대괄호 없음)로, SHA-256 해시의 앞 8자리만 사용하여 원본 복원 불가. 하이드레이션은 대괄호 포함/미포함 형식 모두 지원하여 하위 호환성을 보장한다. 24시간 TTL로 세션 자동 만료.

\subsection{ADR 기록}

프로젝트의 주요 보안 관련 의사결정은 Architecture Decision Record(ADR)로 기록되었다:

\begin{itemize}[nosep]
  \item \textbf{ADR-001}: 키/관리 엔드포인트 인증 없음 — 로컬 단일 사용자 환경에서 계정 관리 비용이 보안 이점보다 큼
  \item \textbf{ADR-002}: 관리 UI용 공개 엔드포인트 — 브라우저 기반 독립 관리 UI가 별도 인증 없이 동작해야 함
  \item \textbf{ADR-003}: SQLite를 기본 데이터베이스로 사용 — PostgreSQL 미실행 환경에서 즉시 동작
  \item \textbf{ADR-004}: 구 MCP 테스트 파일 보존 — 참고용으로 보존, 현행 테스트는 별도 파일에서 수행
\end{itemize}

% ===========================================================================
\section{스마트닝: Two-Phase Extraction}
% ===========================================================================

\subsection{Phase 1 (Extract)}

SLM이 사용자 텍스트를 분석하여 민감 정보를 탐지한다. v1 프롬프트는 패턴 매칭 기반으로, 키워드 단서(숫자 패턴, 회사명, 예산 표현 등)를 열거한다:

\begin{verbatim}
## Step 1: Scan for sensitive spans
- Numbers matching ID formats
- Company names + project/process mentions
- Numbers with "억원", "만원" (budget)
- Phrases about decisions ("결정", "채택", "검토 중")
\end{verbatim}

\subsection{Phase 2 (Critic)}

다른 SLM이 Phase 1 결과를 검토하고 누락된 레코드를 보완한다. 병합 시 다음 검증을 수행한다:

\begin{itemize}[nosep]
  \item 동일한 span + category 조합의 중복 제거
  \item 원본 텍스트에 존재하지 않는 span(할루시네이션) 자동 제거
  \item Phase 2가 Phase 1의 레코드를 수정한 경우 업데이트
\end{itemize}

\subsection{v2 프롬프트: 맥락적 추론 질문}

v2 프롬프트는 구체적 키워드 단서를 유지하되, ``왜 이게 민감한가?''라는 질문을 먼저 던지도록 설계되었다:

\begin{verbatim}
## Step 1c. Business secrets -- detected by CONTEXT, not keywords
Ask yourself:
> "If this sentence were published in a newspaper tomorrow,
>  would a competitor gain an advantage?"

## Step 1d. Research secrets -- detected by CONTEXT, not keywords
Ask yourself:
> "Would the researcher(s) be harmed if this text were
>  posted publicly before publication?"
\end{verbatim}

\subsection{측정 개선치}

v1에서 v2로의 프롬프트 변경만으로 추가 비용 없이 6개 모델에서 사업/연구 기밀 탐지율이 획기적으로 개선되었다:

\begin{table}[h]
\centering
\begin{tabular}{@{}llllll@{}}
\toprule
모델 & \$/1M tok & v1 사업 & v1 연구 & v2 사업 & v2 연구 \\
\midrule
\texttt{gemini-3.1-flash-lite} & \$0.25 & 5/5 & 5/5 & \checkmark & \checkmark \\
\texttt{ministral-3b-2512} & \$0.10 & 0/5 & 0/5 & \checkmark & \checkmark \\
\texttt{deepseek-v4-flash} & \$0.10 & 2/5 & 1/5 & \checkmark & \checkmark \\
\texttt{gemma-4-26b-a4b-it} & \$0.06 & 0/5 & 0/5 & \checkmark & \checkmark \\
\texttt{granite-4.1-8b} & \$0.05 & 0/5 & 0/5 & \checkmark & \checkmark \\
\texttt{qwen3.5-9b} & \$0.04 & 1/5 & 0/5 & \checkmark & \checkmark \\
\texttt{gemini-3.5-flash} & \$1.50 & 0/5 & 0/5 & \checkmark & \checkmark \\
\bottomrule
\end{tabular}
\caption{v1$\rightarrow$v2 사업/연구 기밀 탐지율 비교 (N=5$\rightarrow$단일 패스)}
\end{table}

핵심 발견: 프롬프트 변경만으로 \texttt{ministral-3b-2512}(\$0.10)가 \texttt{gemini-3.1-flash-lite}(\$0.25)와 동등한 7/7 완전 탐지를 달성하여 비용 60\% 절감이 가능하다.

% ===========================================================================
\section{사용 로그 분석 (Usage Log Analysis)}
% ===========================================================================

\subsection{Ground Truth 테스트 데이터}

Privacy Router의 Ground Truth 데이터셋은 총 27개 테스트 케이스로 구성된다:

\begin{table}[h]
\centering
\begin{tabular}{@{}lr@{}}
\toprule
항목 & 수 \\
\midrule
총 테스트 케이스 & 27 \\
\midrule
\multicolumn{2}{l}{\textbf{정책 분포}} \\
\quad \texttt{selective\_mask} (마스킹 후 전송) & 9 \\
\quad \texttt{block} (로컬 처리) & 15 \\
\quad \texttt{allow} (직접 전송) & 3 \\
\midrule
\multicolumn{2}{l}{\textbf{카테고리 분포}} \\
\quad Primary (주요 민감) & 14 \\
\quad Edge (경계 케이스) & 7 \\
\quad Non-sensitive (비민감) & 3 \\
\midrule
\multicolumn{2}{l}{\textbf{난이도 분포}} \\
\quad Easy & 8 \\
\quad Medium & 5 \\
\quad Hard & 11 \\
\bottomrule
\end{tabular}
\caption{Ground Truth 데이터셋 통계}
\end{table}

\subsection{평가 지표}

\texttt{eval\_all.py}가 측정하는 3가지 지표:

\begin{itemize}[nosep]
  \item \textbf{target\_ok}: 민감 여부 판정이 Ground Truth와 일치하는가 (binary)
  \item \textbf{context\_ok}: 정책 결정(allow/mask/block)이 Ground Truth와 일치하는가 (binary)
  \item \textbf{ok}: target\_ok AND context\_ok (최종 통과 기준)
\end{itemize}

라벨링은 프로젝트 팀원 2명이 독립적으로 수행하고, 불일치 시 토론으로 합의했다. 총 17개 케이스에 대해 Cohen's $\kappa > 0.85$를 달성했다.

\subsection{7일 평가 요약}

2단계 평가(Phase 1: 모델 선정, Phase 2: 프롬프트 최적화)를 통해 확인한 주요 결과:

\begin{itemize}[nosep]
  \item Phase 1 (v1, N=5): 9개 모델 평가. \texttt{gemini-3.1-flash-lite}만 사업/연구 기밀 5/5 탐지
  \item Phase 2 (v2, 단일 패스): 프롬프트 변경으로 6개 모델이 7/7 완전 탐지 달성
  \item 로컬 vLLM BF16 평가: Gemma 4 계열에서 E4B(4B) = 26B MoE와 동일한 76.5\% context accuracy 확인
  \item 비선형 스케일링: 2B(41.2\%) $\rightarrow$ 4B(76.5\%) $\rightarrow$ 12B(55.6\%) $\rightarrow$ 26B MoE(76.5\%)
\end{itemize}

\subsection{런타임 사용 로그 (Hermes Agent)}

Hermes Agent가 Privacy Router를 경유하여 실제 작업을 수행한 17분간(2026-06-17 03:52--04:10 UTC)의 라이브 세션을 캡처하였다. 에이전트는 \texttt{gemini-3.1-flash-lite} 모델을 통해 문서 편집, 연구 쿼리, 코드 생성 등 실제 작업을 수행하였으며, 46개 요청 모두 HTTP 200으로 오류 없이 완료되었다.

\begin{table}[h]
\centering
\begin{tabular}{@{}lr@{}}
\toprule
항목 & 값 \\
\midrule
총 요청 수 & 46 \\
민감 정보 포함 요청 & 31 (67.4\%) \\
안전한 요청 & 15 (32.6\%) \\
\midrule
\multicolumn{2}{l}{\textbf{라우팅 정책}} \\
\quad 외부 API로 전달 & 37 (80.4\%) \\
\quad 로컬 처리로 라우팅 & 9 (19.6\%) \\
\midrule
요청당 평균 민감 레코드 수 & 6.1 \\
요청당 최대 민감 레코드 수 & 16 \\
오류율 & 0\% \\
\bottomrule
\end{tabular}
\caption{런타임 사용 통계 (46 요청, 17분)}
\end{table}

주요 관찰 결과:

\begin{itemize}[nosep]
  \item \textbf{요청의 67\%에 민감 정보 포함}: 실제 에이전트 작업의 대부분이 PII 또는 기밀 맥락을 포함하는 프롬프트로, 자동 프라이버시 보호의 필요성을 확인.
  \item \textbf{마스킹 후 80\% 외부 전달}: 대부분의 민감 요청은 안전하게 마스킹하여 전달할 수 있으며, \texttt{is\_essential} 휴리스틱이 마스킹 시 쿼리 의미 보존 여부를 정확히 판단함을 입증.
  \item \textbf{20\% 로컬 처리}: 민감 정보가 쿼리의 대상인 요청(예: ``이 주민등록번호를 포함한 이메일 작성'')은 외부 전송이 올바르게 차단됨.
  \item \textbf{오류율 0\%}: 46개 요청 모두 HTTP 200으로 성공적으로 완료되어 핵심 파이프라인의 프로덕션 안정성 확인.
\end{itemize}

\subsection{전수조사 결과}

프로젝트 전수조사(95 Python + 11 HTML + 22 Config 파일)에서 발견된 문제:

\begin{table}[h]
\centering
\begin{tabular}{@{}lrrr@{}}
\toprule
심각도 & 원래 수 & 수정됨 & 남은 수 \\
\midrule
CRITICAL & 6 & 4 & 0 (즉시 수정 완료) \\
HIGH & 8 & 3 & 0 (즉시 수정 완료) \\
MEDIUM & 9 & 0 & 9 (권장) \\
LOW & 8 & 0 & 8 (선택) \\
\bottomrule
\end{tabular}
\caption{전수조사 결과 요약}
\end{table}

% ===========================================================================
\section{Big-Tech 대비 차별화}
% ===========================================================================

Privacy Router는 Big-Tech 프라이버시 솔루션 대비 7개 차원에서 차별화된다:

\begin{table}[h]
\centering
\begin{tabular}{@{}p{2.5cm}p{4cm}p{4cm}@{}}
\toprule
차원 & Big-Tech & Privacy Router \\
\midrule
1. 데이터 주권 & 클라우드 종속, 데이터 외부 전송 & 로컬 실행, 데이터 외부 미전송 \\
2. 비용 & 월 \$50+ (Azure AI) & 월 \$0.075 \\
3. 다국어 지원 & 영어 중심, 비영어 PII 약함 & 주민등록번호·현지 금액 등 특화 \\
4. 컨텍스트 탐지 & 패턴 매칭만 가능 & 맥락적 추론으로 사업/연구 기밀 탐지 \\
5. 아키텍처 투명성 & 폐쇄형 API & 오픈소스, 전체 파이프라인 가시화 \\
6. 통합 방식 & 전용 SDK 필요 & OpenAI Compatible API로 투명 통합 \\
7. 커스터마이즈 & 제한적 설정 & 에이전트별 모델 지정, 프롬프트 조정 가능 \\
\bottomrule
\end{tabular}
\caption{7차원 비교: Big-Tech vs Privacy Router}
\end{table}

\subsection{통합 실험 결과}

Hermes Agent와 OpenClaw Agent를 Privacy Router API 경유 방식으로 연결한 실험에서:

\begin{itemize}[nosep]
  \item \textbf{Hermes}: Telegram 봇 연결 성공, classify에서 주민등록번호(0.99)·예산(0.99) 정확 탐지, \texttt{route\_to_local} 정책 정상 동작
  \item \textbf{OpenClaw}: Privacy Router를 OpenAI 호환 프로바이더로 연결 성공, 일반 대화는 외부 API 직접 전송, 민감 정보 포함 시 로컬 처리로 자동 전환
\end{itemize}

% ===========================================================================
\section{결론 및 향후 과제}
% ===========================================================================

\subsection{주요 발견}

\begin{enumerate}[nosep]
  \item 컨텍스트 기반 민감 정보 탐지는 모델 크기가 아니라 \textbf{프롬프트 설계}에 의해 결정된다.
  \item 맥락적 추론 질문 2개(``신문에 실리면?'', ``출판 전에 공개되면?'')만으로 모든 테스트 모델의 사업/연구 기밀 탐지율을 0\%$\rightarrow$100\%로 개선했다.
  \item \texttt{ministral-3b-2512}(\$0.10)는 v2 프롬프트 기준 \texttt{gemini-3.1-flash-lite}(\$0.25)와 동등한 성능을 보여, 비용 60\% 절감이 가능하다.
  \item Two-Tier 구성(Extractor: ministral-3b + Judge: Rule-based)으로 월 \$0.075 수준의 운영 비용을 달성했다.
\end{enumerate}

\subsection{한계}

\begin{itemize}[nosep]
  \item v2 단일 패스: 통계적 신뢰도가 낮아 신뢰 구간 추정 불가
  \item 대명사 참조 미처리: ``그것'', ``이 연구'' 등 이전 컨텍스트 참조 표현 탐지 불가
  \item 한국어+영어 테스트 완료: 중국어, 일본어 PII 패턴 미검증
  \item 적대적 테스트 없음: 프롬프트 인젝션, 우회 시도 등 미평가
  \item 레이턴시: SLM 추론 기반 $\sim$2--4초 (경량 token classifier로 대체 계획)
\end{itemize}

\subsection{향후 과제}

\begin{enumerate}[nosep]
  \item v2 프롬프트로 N=5 재평가 $\rightarrow$ 신뢰 구간 확보
  \item Judge 프롬프트에도 동일한 맥락적 추론 기법 적용
  \item 경량 token classifier 기반 Extractor 프로토타입 $\rightarrow$ sub-second 레이턴시 달성
  \item 다국어 확장: 중국어, 일본어 PII 패턴 추가
  \item 적대적 테스트 체계 구축
  \item 대화형 모드(Multi-turn): \texttt{needs\_input} 응답 + 사용자 선택 처리
\end{enumerate}

\end{document}
